\section{犯罪嫌疑人、被告人逃匿、死亡案件违法所得的没收程序}
\begin{enumerate}
  \item \zd{人没了钱还在}，不会作出宣判。
  \item 适用条件。
        \begin{enumerate}
          \item 逃匿：\zd{贪恐重大逃 1 年}。
          \item 死亡：\zd{死了直接收赃钱}。没有罪名、情节限制，只要有赃款赃物。
        \end{enumerate}
  \item 没收对象：赃款赃物及其擎息、违禁品、本人的犯罪工具。
  \item 启动。公安监委\zd{没收违法所得意见书} \yjt 检察院 \yjt 法院。\zd{门当户对}；\zd{捕诉部门}审查，30 + 15。
  \item \zd{程序转换}。
        \begin{enumerate}
          \item 没收程序中 \yjt 人来了 \yjt 终止审查活动 \yjt 转公诉程序 \yjt 退到公安。人没了，没转公。
          \item 公诉程序中 \yjt 人没了 \yjt 不起诉 \yjt 向法院去申请没收违法所得。人没了，公转没。
        \end{enumerate}
  \item 法院对没收程序的审理。
        \begin{enumerate}
          \item 最低中院，合议庭。
          \item 可以查扣冻。防止转移财产 \hz 销毁财产。
          \item 15 天公告，公告期 6 个月，超出公告期也是可能准许的。
          \item 申请参诉的主体：近亲属和利害关系人（除了近亲属之外主张财产所有权的）。
          \item 可以不开庭的，\zd{利害关系人申请参诉的一定要开庭}。
          \item 裁定没收/驳回 \yjt 上诉抗诉期只有 5 天 \yjt 二审程序：
                \begin{enumerate}
                  \item 没错就维持。
                  \item 有错那就应当改。事实不清 \yjt 查清改 \hz 发回重审；程序违法 \yjt 撤销原判裁定，发回重审。
                \end{enumerate}
        \end{enumerate}
  \item 法院审理程序中的转换。\zd{人来了，没转公；人没了，公转没收或缺席}。法院不告不理，永远不主动。
        \begin{enumerate}
          \item 到案 \yjt 终止审理。检察院向 \zd{原受理申请的法院}提起公诉。
          \item 脱逃/死亡 \yjt 裁定中/终止审理 \yjt 满足没收/缺席 \yjt 提起没收/缺席程序。
        \end{enumerate}
  \item 审限参照一审 \hz 二审。
  \item 裁定生效后的救济。
        \begin{enumerate}
          \item 到案，被提起公诉。原裁定正确 \yjt 予以维持；错误 \yjt 应当撤销原裁定，在判决中对涉案财产一并处理。
          \item 没有到案。按照审判监督程序予以纠正。
        \end{enumerate}
\end{enumerate}
